% =====================================================================
%  Lecture 06 -- E-UTRAN: the LTE Radio Access Network   (English edition)
% =====================================================================
\documentclass[11pt,a4paper]{article}
\providecommand\ACNroot{../..}
\input{\ACNroot/templates/lecture-en}

\ACNset{number}{06}
\ACNset{title}{E-UTRAN: the LTE Radio Access Network}
\ACNset{shorttitle}{E-UTRAN}
\ACNset{titleother}{E-UTRAN: شبکهٔ دسترسی رادیویی LTE}
\ACNset{part}{Part II -- 4G access and core (Lectures 5--13)}
\ACNset{time}{7--9 hours (3.5 h reading, 4--5 h exercises)}
\ACNset{level}{Intermediate}
\ACNset{prereq}{Lecture 5; the OSI layering; basic ARQ}
\ACNset{gen}{4}

\begin{document}
\ACNcover

% ---------------------------------------------------------------------
\section{Learning objectives}
After studying this lecture you should be able to:
\begin{objectives}
  \item \textbf{List} the functions of the eNB and explain why the flat E-UTRAN could absorb
        the RNC.
  \item \textbf{Draw} the LTE control-plane and user-plane protocol stacks and state the role of
        PHY, MAC, RLC, PDCP and RRC.
  \item \textbf{Map} logical channels to transport and physical channels in both directions.
  \item \textbf{Explain} the S1 and X2 interfaces, their protocol stacks and main procedures.
  \item \textbf{Describe} the RRC state machine, DRX and the NR RRC\_INACTIVE extension, and
        \textbf{quantify} the power/latency trade-off of DRX.
  \item \textbf{Size} PDCP sequence numbers and RLC windows for a target data rate.
\end{objectives}

% ---------------------------------------------------------------------
\section{Prerequisite knowledge}
This lecture zooms into the left half of the EPS architecture of Lecture~5. You need the
idea of protocol layering (each layer offers a service to the one above and adds a header),
the notion of automatic repeat request (ARQ) with sequence numbers and windows, and the
EPS bearer concept.

\paragraph{Key terms} The vocabulary of this lecture with its Persian equivalents:
\begin{keyterms}
E-UTRAN & \fa{شبکهٔ دسترسی رادیویی LTE} & The LTE RAN: a mesh of eNBs, no controller. \\
RRC & \fa{کنترل منابع رادیویی} & Layer-3 control between UE and eNB: connection, measurements, security. \\
PDCP & \fa{پروتکل همگرایی دادهٔ بسته‌ای} & Header compression, ciphering, integrity, in-order delivery. \\
RLC & \fa{کنترل پیوند رادیویی} & Segmentation and ARQ in TM, UM or AM mode. \\
MAC & \fa{کنترل دسترسی به رسانه} & Multiplexing, scheduling, HARQ, random access. \\
Logical channel & \fa{کانال منطقی} & \emph{What} is carried (control or traffic). \\
Transport channel & \fa{کانال انتقال} & \emph{How} it is carried (format, coding) between MAC and PHY. \\
Physical channel & \fa{کانال فیزیکی} & The actual time--frequency resources. \\
S1AP / X2AP & \fa{پروتکل‌های کاربردی S1 و X2} & Signalling eNB--MME and eNB--eNB over SCTP. \\
DRX & \fa{دریافت ناپیوسته} & Sleep cycles that save battery at the cost of latency. \\
\end{keyterms}

% ---------------------------------------------------------------------
\section{Introduction and motivation}
In UMTS the radio protocols were split across two nodes: the NodeB handled the physical layer,
the RNC everything above it. Every scheduling decision, every retransmission of RLC and every
handover therefore crossed the Iub transport network, adding delay and making the RNC a
bottleneck. HSPA already moved fast scheduling and HARQ to the NodeB. LTE completed the move:
the eNB terminates \emph{all} radio protocols, so the round trip between a radio event and a
decision stays inside one box.

This lecture opens that box. We follow user data and signalling down the protocol stack, see
how the different kinds of information are organised into channels, and study the two network
interfaces of the eNB: S1 towards the core and X2 towards its neighbours. We end with the RRC
state machine, which decides how much the UE and the network must keep in memory and how fast
an idle phone can start sending data.

% ---------------------------------------------------------------------
\section{Fundamental concepts}

\subsection{Why no RNC?}
The RNC existed for three reasons: soft-handover combining, central radio resource management
of many NodeBs, and termination of protocols that were too complex for 1990s base stations.
LTE removed each reason: OFDMA uses hard handover, so no combining is needed; scheduling every
1\,ms must be done next to the radio; and processing power had become cheap. Inter-cell
coordination that the RNC used to provide is replaced by direct eNB-to-eNB signalling over X2
and by self-organising network functions.

\subsection{Layering and the bearer service}
Each layer of the radio stack offers a well-defined service upwards:
\begin{itemize}
  \item \textbf{PHY} offers \emph{transport channels} to MAC.
  \item \textbf{MAC} offers \emph{logical channels} to RLC.
  \item \textbf{RLC} offers \emph{RLC channels} to PDCP.
  \item \textbf{PDCP} offers \emph{radio bearers} to IP (data radio bearers, DRBs) and to RRC
        (signalling radio bearers, SRBs).
\end{itemize}
A data radio bearer is the radio segment of an EPS bearer (Lecture~5); SRB0 carries RRC messages
before security is active, SRB1 and SRB2 carry RRC and NAS messages afterwards.

% ---------------------------------------------------------------------
\section{Detailed technical explanation}

\subsection{eNB functions}
The eNB performs \cite{ts36300}:
\begin{itemize}
  \item \textbf{Radio resource management}: radio bearer control, radio admission control,
        connection mobility control (handover), dynamic allocation of resources to UEs in uplink
        and downlink (scheduling), inter-cell interference coordination (ICIC), load balancing.
  \item \textbf{Header compression and encryption} of user data (PDCP).
  \item \textbf{MME selection} at attach (NAS node selection function) and routing of user data
        towards the S-GW.
  \item \textbf{Paging} and \textbf{system information broadcast}, originating from the MME and
        O\&M respectively.
  \item \textbf{Measurement configuration} for mobility and scheduling.
\end{itemize}

\subsection{The protocol layers}
\paragraph{Physical layer (PHY)} OFDMA in the downlink and SC-FDMA in the uplink with
\qty{15}{\kilo\hertz} subcarriers; turbo coding for data; QPSK, 16QAM and 64QAM (256QAM in later
releases); up to $4\times4$ (later $8\times8$) MIMO; HARQ soft combining; cell search with
primary and secondary synchronisation signals (PSS/SSS). The resource grid and its numerology are
the subject of Lecture~7.

\paragraph{MAC} Maps logical channels to transport channels, multiplexes several logical channels
into one transport block, performs the \emph{scheduling} decision every TTI, runs HARQ (8 parallel
stop-and-wait processes in FDD), handles logical-channel priorities, random access (Lecture~8),
timing advance and buffer status and power headroom reports \cite{ts36321}.

\paragraph{RLC} Has three modes \cite{ts36322}:
\begin{itemize}
  \item \textbf{Transparent mode (TM)}: no header, no segmentation -- used for broadcast, paging
        and SRB0.
  \item \textbf{Unacknowledged mode (UM)}: segmentation/concatenation and reordering, no
        retransmission -- used for delay-sensitive traffic such as VoLTE.
  \item \textbf{Acknowledged mode (AM)}: adds ARQ with status reports -- used for TCP traffic and
        signalling, where reliability matters more than delay.
\end{itemize}
RLC adapts IP packets to the transport-block size chosen by the scheduler: it may cut a packet
into segments or pack several into one PDU.

\paragraph{PDCP} \cite{ts36323} performs robust header compression (ROHC) of IP/UDP/RTP headers,
ciphering of user and control data, integrity protection of control data (SRBs), and in-order,
duplicate-free delivery across handover (using the SN status transfer of Lecture~4). Each PDCP
PDU carries a sequence number (SN); together with a hyper-frame number (HFN) kept at both ends it
forms the 32-bit \emph{COUNT}, an input to the ciphering algorithm that must never repeat under the
same key.

\paragraph{RRC} \cite{ts36331} broadcasts system information (MIB on PBCH, SIBs on PDSCH), manages
the RRC connection (setup, reconfiguration, re-establishment, release), configures measurements
and receives reports, activates access-stratum security, and transfers NAS messages.

\subsection{Channels}
\Cref{fig:channels} shows the main mapping. The logical channels are the broadcast (BCCH), paging
(PCCH), common control (CCCH, before an RRC connection exists), dedicated control (DCCH) and
dedicated traffic (DTCH) channels. Transport channels group them by transmission format: the
downlink and uplink shared channels (DL-SCH, UL-SCH) carry almost everything. Physical channels
include the shared data channels PDSCH and PUSCH, control channels without transport channels
(PDCCH for scheduling assignments, PCFICH for the size of the control region, PHICH for uplink
HARQ acknowledgements, PUCCH for uplink control) and reference signals (CRS, DMRS, SRS) used for
measurements and demodulation.

\begin{acnfigure}
  \resizebox{\linewidth}{!}{\acnfig{lecture-06/channel-mapping}}
  \caption{Main LTE channel mapping. Almost all logical channels converge on the shared channels;
  PDCCH, PHICH, PCFICH and PUCCH carry layer-1/2 control without a transport channel.}
  \label{fig:channels}
\end{acnfigure}

\subsection{The S1 interface}
S1 connects each eNB to one or more MMEs (S1-MME) and S-GWs (S1-U) -- a many-to-many relation
called \emph{S1-flex} that allows MME pooling (Lecture~12).
\begin{itemize}
  \item \textbf{S1-MME}: S1AP over SCTP over IP \cite{ts36413}. SCTP provides reliable, message-oriented
        transport with multi-homing and multiple streams, avoiding head-of-line blocking between
        UEs.
  \item \textbf{S1-U}: GTP-U over UDP over IP; one tunnel per E-RAB (the S1 part of an EPS bearer).
\end{itemize}
Main S1AP procedures: \emph{S1 Setup} (eNB and MME exchange identities and served tracking areas),
\emph{Initial UE Message} (first NAS message of a UE), \emph{Initial Context Setup} (MME installs
the UE context: security key, bearers, capabilities), \emph{E-RAB Setup/Modify/Release},
\emph{Paging}, \emph{Handover Preparation/Resource Allocation} (S1 handover), \emph{Path Switch},
\emph{UE Context Release} and \emph{Downlink/Uplink NAS Transport}.

\subsection{The X2 interface}
X2 connects neighbouring eNBs \cite{ts36423}: X2AP over SCTP for control, GTP-U for forwarding user
data during handover. Main procedures: \emph{X2 Setup} (exchange of served cells, enabling automatic
neighbour relations), \emph{Handover Preparation}, \emph{SN Status Transfer}, \emph{UE Context
Release}, \emph{Load Indication} (interference information for ICIC) and \emph{Resource Status
Reporting} (load for mobility load balancing). X2 is optional: without it, handovers go through
S1 and the MME.

\subsection{RRC states and DRX}
LTE has two RRC states (\cref{fig:rrc}):
\begin{itemize}
  \item \textbf{RRC\_IDLE}: no RRC connection, no context in the eNB; the UE performs cell
        reselection, reads system information and monitors paging once per paging cycle (idle DRX);
        the MME knows it at tracking-area level (ECM-IDLE).
  \item \textbf{RRC\_CONNECTED}: the eNB holds the UE context; the UE can be scheduled, reports
        measurements and is handed over by the network; to save power it may use connected-mode
        DRX, waking for an \emph{on-duration} once per DRX cycle.
\end{itemize}
Moving from IDLE to CONNECTED requires random access, RRC connection setup, a NAS service request,
security activation and bearer set-up -- tens of milliseconds and a dozen messages. NR therefore
adds \textbf{RRC\_INACTIVE} \cite{ts38331}: the UE and the gNB keep the context (security,
bearers), the core still sees the UE as connected, and resumption needs only a short exchange.

\begin{acnfigure}
  \acnfig{lecture-06/rrc-states}
  \caption{RRC state machine. LTE has IDLE and CONNECTED; NR adds INACTIVE, in which the context is
  suspended rather than released.}
  \label{fig:rrc}
\end{acnfigure}

% ---------------------------------------------------------------------
\section{Mathematical models}

\subsection{Header compression gain}\label{sec:rohc}
\paragraph{Problem} How much radio capacity does ROHC save for VoLTE?
\paragraph{Model} A speech frame of $L_\mathrm{p}$ bytes is carried with $H$ bytes of RTP/UDP/IP
headers; ROHC reduces $H$ to a compressed header of $H_\mathrm{c}$ bytes in steady state. The
radio efficiency of the packet is
\begin{equation}\label{eq:rohc}
  \eta = \frac{L_\mathrm{p}}{L_\mathrm{p} + H}, \qquad
  G = \frac{L_\mathrm{p}+H}{L_\mathrm{p}+H_\mathrm{c}} \ \text{(capacity gain)}.
\end{equation}
\paragraph{Example} $L_\mathrm{p} = 32$\,bytes, $H = 40$\,bytes (IPv4 20 + UDP 8 + RTP 12),
$H_\mathrm{c} = 3$\,bytes: $\eta$ rises from $44\%$ to $91\%$ and $G = 72/35 = 2.06$. With IPv6
($H = 60$) the gain is $92/35 = 2.6$. ROHC roughly doubles the number of VoLTE users a cell can
carry; it is not applied on S1-U, where the headers are carried in full (Lecture~5).
\paragraph{Limits} The steady-state compressed size is reached only after a context is established;
packet loss can desynchronise the decompressor, so ROHC profiles trade compression against
robustness.

\subsection{Sequence-number space and wrap-around}\label{sec:sn}
\paragraph{Problem} How many PDCP SN bits are needed at a given rate?
\paragraph{Model} A window-based protocol with an $n$-bit SN can have at most $2^{n-1}$ PDUs
outstanding (half the space), otherwise old and new PDUs become ambiguous. With a PDU rate
$\lambda_\mathrm{p}$ (PDUs per second), the window covers the time
\begin{equation}\label{eq:snwin}
  T_{\mathrm{win}} = \frac{2^{\,n-1}}{\lambda_\mathrm{p}}, \qquad
  \lambda_\mathrm{p} = \frac{R}{8L},
\end{equation}
for a bit rate $R$ and PDUs of $L$ bytes. $T_{\mathrm{win}}$ must exceed the longest reordering
or retransmission delay, typically including HARQ, RLC ARQ and handover data forwarding.
\paragraph{Example} $R = \qty{300}{\mega\bit\per\second}$, $L = 1500$\,bytes:
$\lambda_\mathrm{p} = 25\,000$\,PDU/s. A 12-bit SN gives $T_{\mathrm{win}} = 2048/25\,000 = 82$\,ms
-- uncomfortably short during handover. This is why LTE added a 15-bit (655\,ms) and NR an 18-bit
PDCP SN.

\subsection{RLC acknowledged-mode throughput}\label{sec:rlcwin}
For a window of $W$ PDUs of $L$ bytes and a status round-trip time $T_{\mathrm{RTT}}$, a
window-limited sender cannot exceed
\begin{equation}\label{eq:rlc}
  R_{\max} = \frac{8\,W\,L}{T_{\mathrm{RTT}}} .
\end{equation}
With the LTE RLC AM 10-bit SN, $W = 512$; for $L = 1500$\,bytes and $T_{\mathrm{RTT}} = 20$\,ms,
$R_{\max} = 512\times12\,000/0.02 = \qty{307}{\mega\bit\per\second}$: enough for Release~8, the
reason for longer SNs in later releases with higher rates (\cref{exr:rlc}).

\subsection{DRX: power versus latency}\label{sec:drx}
\paragraph{Problem} How long should a connected UE sleep?
\paragraph{Model} DRX cycle $T_{\mathrm{DRX}}$ with on-duration $T_{\mathrm{on}}$; power
$P_{\mathrm{on}}$ when awake and $P_{\mathrm{sleep}}$ when asleep; downlink packets arrive at random
instants (no traffic during the sleep period otherwise).
\paragraph{Derivation} With duty cycle $d = T_{\mathrm{on}}/T_{\mathrm{DRX}}$,
\begin{equation}\label{eq:drx}
  \bar P = d\,P_{\mathrm{on}} + (1-d)\,P_{\mathrm{sleep}}, \qquad
  \bar D_{\mathrm{DRX}} \approx \frac{T_{\mathrm{DRX}}}{2},
\end{equation}
because a packet arriving at a uniformly random time waits on average half a cycle for the next
on-duration.
\paragraph{Example and interpretation} $T_{\mathrm{on}} = 10$\,ms, $P_{\mathrm{on}} = 500$\,mW,
$P_{\mathrm{sleep}} = 10$\,mW (illustrative): a 40\,ms cycle gives 132.5\,mW and 20\,ms extra delay;
a 320\,ms cycle 25.3\,mW and 160\,ms; 2560\,ms gives 11.9\,mW but 1.28\,s (\cref{fig:drx}). Power falls
like $1/T_{\mathrm{DRX}}$ until $P_{\mathrm{sleep}}$ dominates, while latency grows linearly: beyond a
few hundred milliseconds longer cycles buy little energy at a large latency cost. Real devices
also stay awake for an inactivity timer after each packet, which the model ignores.

\begin{acnfigure}
  \acnfig{lecture-06/drx-power}
  \caption{Average power (left axis) and mean added latency (right axis) versus DRX cycle for
  $T_{\mathrm{on}} = 10$\,ms, from \cref{eq:drx}.}
  \label{fig:drx}
\end{acnfigure}

% ---------------------------------------------------------------------
\section{Network architecture and diagrams}
\Cref{fig:stacks} summarises the radio stacks. On the control plane, RRC terminates in the eNB
while NAS passes through it transparently to the MME (carried in RRC towards the UE and in S1AP
towards the MME). On the user plane, PDCP, RLC and MAC terminate in the eNB, which relays the IP
packets into GTP-U tunnels on S1-U.

\begin{acnfigure}
  \resizebox{\linewidth}{!}{\acnfig{lecture-06/protocol-stacks}}
  \caption{LTE radio protocol stacks: control plane (left) and user plane (right).}
  \label{fig:stacks}
\end{acnfigure}

% ---------------------------------------------------------------------
\section{Step-by-step procedures}
\begin{procedurebox}{From RRC\_IDLE to data transfer (mobile-originated)}\label{proc:idle2conn}
\begin{steps}
  \item \textbf{Random access.} The UE sends a preamble, receives a random access response with a
        timing advance and grant (Lecture~8).
  \item \textbf{RRC connection request.} On SRB0 it sends \msg{RRC Connection Request} (identity and
        cause); the eNB answers \msg{RRC Connection Setup}, creating SRB1.
  \item \textbf{NAS transfer.} The UE sends \msg{RRC Connection Setup Complete} containing the NAS
        \msg{Service Request}; the eNB forwards it in S1AP \msg{Initial UE Message} to the MME
        selected from the UE's identity.
  \item \textbf{Context set-up.} The MME answers S1AP \msg{Initial Context Setup Request} with the
        security key $K_{\mathrm{eNB}}$, the E-RABs and the S-GW tunnel endpoints.
  \item \textbf{AS security.} The eNB sends \msg{Security Mode Command}; the UE derives the AS keys
        and answers \msg{Security Mode Complete}. PDCP now ciphers and integrity-protects.
  \item \textbf{Bearer configuration.} The eNB sends \msg{RRC Connection Reconfiguration} creating SRB2
        and the DRBs; after \msg{Reconfiguration Complete} it answers \msg{Initial Context Setup Response}
        with its own S1-U tunnel endpoints. Uplink and downlink data flow.
\end{steps}
\end{procedurebox}

% ---------------------------------------------------------------------
\section{Worked examples}

\begin{examplebox}{Protocol overhead of a TCP segment}\label{ex:overhead}
\textbf{Problem.} A 1460-byte TCP payload is sent over LTE with IPv4. Count the header bytes added
on the radio by PDCP (2 bytes, 12-bit SN), RLC AM (2 bytes, no segmentation) and MAC (2 bytes
subheader), and on S1-U.

\textbf{Solution.} IP packet: $1460 + 20 + 20 = 1500$ bytes. Radio: $1500 + 2 + 2 + 2 = 1506$ bytes
(0.4\% overhead; ROHC is usually not used for TCP). S1-U: $1500 + 36 = 1536$ bytes (2.4\%).

\textbf{Interpretation.} For large packets layer-2 overhead is negligible; the real costs are the
reference signals and control channels of the PHY (Lecture~5) and the retransmissions.
\end{examplebox}

\begin{examplebox}{PDCP window during handover}\label{ex:pdcp}
\textbf{Problem.} A UE receives 150\,Mbit/s in 1500-byte packets. During handover, forwarded
packets may arrive up to 60\,ms late. Is a 12-bit PDCP SN sufficient?

\textbf{Solution.} $\lambda_\mathrm{p} = 150\times10^6/12\,000 = 12\,500$\,PDU/s. \Cref{eq:snwin}:
$T_{\mathrm{win}} = 2048/12\,500 = 164$\,ms $> 60$\,ms. Yes, with margin; at 300\,Mbit/s (82\,ms) the
margin would be thin.
\end{examplebox}

\begin{examplebox}{Choosing a DRX cycle}\label{ex:drx}
\textbf{Problem.} A messaging application tolerates 200\,ms extra latency on average. Using the
parameters of \cref{sec:drx}, choose the longest standard DRX cycle among 40, 80, 160, 320\,ms and
compute the power saving relative to no DRX ($P_{\mathrm{on}}$ continuously).

\textbf{Solution.} $T_{\mathrm{DRX}}/2 \le 200$\,ms allows 320\,ms (160\,ms latency). $\bar P = 25.3$\,mW
versus 500\,mW: a 95\% reduction.

\textbf{Interpretation.} Most of the saving is already obtained at 160\,ms (40.6\,mW); the operator
can choose per application through the inactivity timer and cycle settings.
\end{examplebox}

% ---------------------------------------------------------------------
\section{Technology comparison}
\begin{table}[htbp]
\centering\small
\caption{Radio access network functions: UTRAN versus E-UTRAN versus NG-RAN.}\label{tab:ran}
\rowcolors{2}{acnmist}{white}
\begin{tabularx}{\linewidth}{@{}l Y Y Y@{}}
\toprule
\rowcolor{white} & UTRAN (3G) & E-UTRAN (4G) & NG-RAN (5G) \\
\midrule
Nodes              & NodeB + RNC          & eNB                       & gNB (optionally CU + DU) \\
RLC/PDCP location  & RNC                  & eNB                       & DU (RLC), CU (PDCP) \\
Scheduler          & RNC (R99), NodeB (HSPA) & eNB, 1\,ms             & gNB-DU, slot/mini-slot \\
Inter-node link    & Iur                  & X2                        & Xn; F1 between CU and DU \\
Core interface     & Iu (RANAP)           & S1 (S1AP)                 & NG (NGAP) \\
RRC states         & Idle, Cell\_DCH/FACH/PCH, URA\_PCH & Idle, Connected & Idle, Inactive, Connected \\
New L2 sublayer    & --                   & --                        & SDAP (QoS-flow mapping) \\
\bottomrule
\end{tabularx}
\end{table}

\paragraph{Trade-off} Putting everything in the eNB minimised latency but made coordination between
cells (interference, joint transmission) harder, since it now requires X2 signalling between peers.
5G partially re-centralises the RAN with the CU/DU split (Lecture~14): the time-critical lower layers
stay at the antenna site, PDCP and RRC can be pooled -- a deliberate compromise between UTRAN's
centralisation and E-UTRAN's flatness.

% ---------------------------------------------------------------------
\section{Key takeaways}
\begin{takeawaybox}
\begin{itemize}
  \item The eNB terminates PHY, MAC, RLC, PDCP and RRC: radio decisions stay in one node.
  \item PDCP: ROHC, ciphering, integrity (SRBs), in-order delivery; COUNT $=$ HFN $\Vert$ SN must never repeat.
  \item RLC: TM, UM (voice), AM (TCP, signalling); MAC: multiplexing, scheduling, HARQ.
  \item Logical $\to$ transport $\to$ physical channels; almost everything uses the shared channels.
  \item S1-MME (S1AP/SCTP) and S1-U (GTP-U/UDP) towards the core; X2 between eNBs.
  \item RRC\_IDLE vs RRC\_CONNECTED (+ RRC\_INACTIVE in NR); DRX trades power ($\propto1/T_{\mathrm{DRX}}$)
        for latency ($T_{\mathrm{DRX}}/2$).
  \item SN space must cover the reordering time: $2^{n-1}/\lambda_\mathrm{p} > T_{\mathrm{reorder}}$.
\end{itemize}
\end{takeawaybox}

% ---------------------------------------------------------------------
\section{Exercises}

\begin{exercisebox}{1}{Which layer?}\label{exr:layer}
For each function name the LTE layer that performs it: HARQ, ARQ, ciphering of user data,
segmentation, header compression, integrity protection of RRC messages, choice of MCS, measurement
configuration.
\end{exercisebox}

\begin{exercisebox}{1}{Why SCTP?}\label{exr:sctp}
Explain why S1AP and X2AP use SCTP rather than TCP or UDP. Give one failure scenario that
multi-homing handles.
\end{exercisebox}

\begin{exercisebox}{2}{RLC window}\label{exr:rlc}
Using \cref{eq:rlc}, find the RLC AM window needed for 1\,Gbit/s with 1500-byte PDUs and a 15\,ms
status round trip, and the minimum number of SN bits.
\end{exercisebox}

\begin{exercisebox}{2}{ROHC and IPv6}\label{exr:rohc}
Repeat \cref{sec:rohc} for AMR-WB at 23.85\,kbit/s (61-byte frames every 20\,ms) with IPv6. What is
the bit rate on the radio with and without ROHC, and the capacity gain?
\end{exercisebox}

\begin{exercisebox}{3}{DRX with an inactivity timer}\label{exr:drxinact}
After each received packet the UE stays awake for an inactivity timer $T_{\mathrm{in}}$. Packets
arrive as a Poisson process with rate $\lambda$. Derive an approximate expression for the fraction of
time awake, assuming $T_{\mathrm{in}} \ll 1/\lambda$, and evaluate it for $\lambda = 0.2$\,s$^{-1}$,
$T_{\mathrm{in}} = 100$\,ms, $T_{\mathrm{on}} = 10$\,ms, $T_{\mathrm{DRX}} = 320$\,ms.
\end{exercisebox}

\begin{exercisebox}{3}{Idle versus connected}\label{exr:idleconn}
Releasing a UE to RRC\_IDLE after inactivity saves eNB memory and UE power, but each new data burst
then costs about 12 signalling messages. If the eNB can handle $10^4$ such transitions per second and
each UE generates a burst every 30\,s, how many UEs can be served purely from the signalling point of
view? Discuss how RRC\_INACTIVE changes the answer.
\end{exercisebox}

\subsection*{Solutions}
\addcontentsline{toc}{subsection}{Solutions}

\begin{solutionbox}{\ref{exr:layer}}
HARQ: MAC. ARQ: RLC (AM). Ciphering of user data: PDCP. Segmentation: RLC. Header compression:
PDCP. Integrity of RRC: PDCP. Choice of MCS: MAC scheduler (with link adaptation input from PHY).
Measurement configuration: RRC.
\end{solutionbox}

\begin{solutionbox}{\ref{exr:sctp}}
SCTP is message-oriented (signalling messages keep their boundaries), reliable and congestion
controlled like TCP, supports multiple streams so that a lost message of one UE does not block
messages of others (no head-of-line blocking across streams), and supports multi-homing: an
association can use several IP paths. If one transport link between eNB and MME fails, the
association continues over the other address without tearing down the S1 interface and every UE
context with it.
\end{solutionbox}

\begin{solutionbox}{\ref{exr:rlc}}
$W \ge R\,T_{\mathrm{RTT}}/(8L) = 10^9\times0.015/12\,000 = 1250$ PDUs. The SN space must be at least
$2W = 2500$, so 12 bits ($4096$) are needed; LTE Release~8's 10 bits are insufficient (hence the
longer RLC sequence numbers of later releases; NR RLC AM uses 12 or 18 bits).
\end{solutionbox}

\begin{solutionbox}{\ref{exr:rohc}}
IPv6/UDP/RTP: $40+8+12 = 60$ bytes. Without ROHC: $61+60 = 121$ bytes per 20\,ms $=$ 48.4\,kbit/s;
with ROHC (3 bytes): 64 bytes $=$ 25.6\,kbit/s. Gain $121/64 = 1.89$. The gain is smaller than for
32-byte frames because the payload is larger.
\end{solutionbox}

\begin{solutionbox}{\ref{exr:drxinact}}
Each packet keeps the UE awake for about $T_{\mathrm{in}}$; with $\lambda T_{\mathrm{in}} \ll 1$ these
periods rarely overlap. During the remaining time DRX applies. Fraction awake
$f \approx \lambda T_{\mathrm{in}} + (1-\lambda T_{\mathrm{in}})\,T_{\mathrm{on}}/T_{\mathrm{DRX}}
= 0.02 + 0.98\times0.03125 = 0.0506$, i.e.\ 5.1\%. With the powers of \cref{sec:drx},
$\bar P \approx 0.0506\times500 + 0.9494\times10 = 34.8$\,mW. The inactivity timer adds about 40\% to
the energy of the pure DRX case for this traffic.
\end{solutionbox}

\begin{solutionbox}{\ref{exr:idleconn}}
Each UE causes $1/30$ transitions per second, so $N = 10^4\times30 = 3\times10^5$ UEs per eNB --
more than any cell carries, but the MME sees the sum over many eNBs, and paging, NAS and S1 messages
load the core too. With RRC\_INACTIVE a resumption needs only about 3 radio messages and no core
signalling (the core sees the UE as connected), so the core load disappears and the radio load drops
by about $4\times$, at the cost of storing contexts in the RAN.
\end{solutionbox}

% ---------------------------------------------------------------------
\section{Further reading}
The E-UTRAN overall description \cite{ts36300} is the best entry point; the layer specifications are
MAC \cite{ts36321}, RLC \cite{ts36322}, PDCP \cite{ts36323} and RRC \cite{ts36331}; the network
interfaces are S1AP \cite{ts36413} and X2AP \cite{ts36423}; SCTP is specified in \cite{rfc9260}.
Sesia et al.\ \cite{sesia2011lte} devote several chapters to the layer-2 protocols.

\printbibliography[heading=subbibliography,title={\ACNstrRefs}]
\end{document}
